\documentclass[12pt]{article}

%\usepackage[cm]{fullpage}
\setlength{\topmargin}{-1.2in}
\setlength{\oddsidemargin}{-0.25in}
\setlength{\evensidemargin}{-0.25in}
\setlength{\textwidth}{7in}
\setlength{\textheight}{10.2in}

\usepackage{array}
\usepackage{enumitem}
\usepackage{amssymb}
\usepackage{amsmath}
\usepackage{epstopdf}
\usepackage[T1]{fontenc}
\usepackage{graphics}
\usepackage{graphicx}
\usepackage{hyperref}
\usepackage{latexsym}
\usepackage{mathrsfs}
\usepackage{multicol}
\usepackage{setspace}
\usepackage{subfigure}
\usepackage{time}
\usepackage{times}

\newcommand{\To}{\Rightarrow}
\newcommand{\del}{\nabla}
\newcommand{\R}{\mathbb{R}}
\newcommand{\C}{\mathcal{C}}
\newcommand{\N}{\mathcal{N}}
\newcommand{\A}{\mathcal{A}}
\renewcommand{\O}{\mathcal{O}}
% \renewcommand{\N}{\mathbb{N}}
\newcommand{\ben}{\begin{enumerate}}
\newcommand{\een}{\end{enumerate}}
\newcommand{\eps}{\varepsilon}
\renewcommand{\epsilon}{\varepsilon}
\newcommand{\p}{\partial}
\newcommand{\<}{\langle}
\renewcommand{\>}{\rangle}
\renewcommand{\subset}{\subseteq}
\newcommand{\cont}{\Rightarrow \Leftarrow}
\newcommand{\back}{\backslash}
\newcommand{\norm}[1]{\|{#1}\|}
\newcommand{\abs}[1]{|{#1}|}
\newcommand{\ip}[1]{\langle{#1}\rangle}
\newcommand{\m}[1]{{\mathbf{#1}}}
\newcommand{\sts}{\m{S}^\textrm{T}\m{S}}
\newcommand{\rank}{\textrm{rank}}
\newcommand{\X}{\mathcal{X}}
\newcommand{\Y}{\mathcal{Y}}
\def\ds{\displaystyle}
\def\d{\partial}
\renewcommand{\arraystretch}{1.5}
\renewcommand{\L}{\mathscr{L}}

%\includegraphics[scale = 0.65, trim = 1in 3.25in 1in 3in]{blank_axes.pdf}


\pagestyle{empty}
\begin{document}
\noindent\rule{7in}{0.4pt}\\
\begin{center}
{\sc Differential Equations\\
ODE versus PDE}
\end{center}
\noindent\rule{7in}{0.4pt}
%%%%%%%%%%%%%%%%%%%%%%%%%%%%%%%%%%%%%%%%%%%%%%%%
\section*{ODE versus PDE, Independent versus Dependent Variables}
\doublespacing
\noindent  The focus of this text is %{\bf ordinary differential equations}
\line(1,0){400}, where the unknown function
is a function of just one variable.  Consider, for example, the DE\\ \\
%\vspace{0.52in}
%	\[ y'' - y' = e^x. \]

\noindent  Here the unknown function is %$y$
\line(1,0){50}, and whatever this unknown function is, we would want
it to be a function such that \\ \vspace{0.25in}
%if we took its second derivative and subtracted its first derivative, we would get
%$e^x.$  

\noindent It makes sense, then, that %$y$ is actually a function of $x$
\line(1,0){300}, and we are\\ \\
%taking derivatives {\it with respect to $x$}
\line(1,0){500}.  \\ \\ So we have
one %{\bf independent variable}, $x$, 
\line(1,0){300}\, and \\ \\
one %{\bf dependent variable}, our unknown function $y$.  
\line(1,0){500}.\\
\noindent Hence, this is an example of an ordinary differential equation.\\

\noindent  It is possible that the unknown function in a DE depends on more than one independent variable.
Suppose we want to describe a function $u$ that gives the temperature at a point $(x, y, z)$ at any time $t$.
Then the value of $u$ depends on four independent variables, $x$, $y$, $z$, and $t$.  If we want to take
the derivative of $u$, then we need to be careful to indicate if we are taking a derivative with respect to $x$
or $y$ or $z$ or $t$, and multiple of these derivatives might occur in the same DE.  The heat equation
	\[ \frac{\partial u}{\partial t} = \alpha \left( \frac{\partial^2 u}{\partial x^2} 
		+ \frac{\partial^2 u}{\partial y^2} 
		+ \frac{\partial^2 u}{\partial z^2}\right), \hspace{1cm} \alpha >0, \]	%margin note:  it's OK if you don't understand the notation here...
describes the relationship between some derivatives of the function $u$ and shows how heat diffuses
over time in a solid medium.  This type of DE is called a %{\bf partial differential equation}, or PDE.  
\line(1,0){300}.\\  This text will
focus exclusively on ODEs, but it is worth knowing that PDEs are common and there exist solution techniques
to solve those, too.

\end{document}